=====================================================================
 BOOK DETAILS PACKAGE
 For manual entry by the publisher.
 All fields below are plain text; copy each into the corresponding
 field of the publisher's submission form.
=====================================================================

---------------------------------------------------------------------
TITLE
---------------------------------------------------------------------
Title:
  The Blindness of the Bounded
%
Subtitle:
  Epistemic Floors, Loop Closure, and a Reference-Free Theory of
  Cellular Health
%
(Alternative, more formal subtitle if the publisher prefers a
 domain-first phrasing:)
  A Mathematical Theory of Bounded Observation, Self-Consistency,
  and Disease in Cellular Systems


---------------------------------------------------------------------
SPINE
---------------------------------------------------------------------
Spine title (shortened for the narrow spine):
  The Blindness of the Bounded
%
Spine author:
  Sachikonye
%
Spine colours:
  Background:  deep slate blue-grey       HEX #2B3A4A   (RGB 43, 58, 74)
  Lettering:   warm off-white / bone       HEX #EDE6D6   (RGB 237, 230, 214)
  Accent rule: muted ochre                 HEX #B8894B   (RGB 184, 137, 75)
%
(Rationale: the slate blue-grey and bone recall the tempera palette of
 Bruegel's field and sky; the ochre accent recalls the earth of the
 ditch and the workshop's azurite-and-ultramarine restraint. Sober,
 scholarly, not decorative.)


---------------------------------------------------------------------
CLASSIFICATION
---------------------------------------------------------------------
Field / subject category:
  Mathematical biology; theoretical systems biology
  (secondary: applied mathematics; biophysics; philosophy of science)
%
BISAC / subject codes (suggested):
  SCI008000  SCIENCE / Life Sciences / Biology
  MAT003000  MATHEMATICS / Applied
  SCI013000  SCIENCE / Chemistry / Physical & Theoretical
  MED000000  MEDICAL / General  (secondary; for the therapeutic chapters)
%
Language:
  English


---------------------------------------------------------------------
KEYWORDS
---------------------------------------------------------------------
(comma-separated, for catalogue/search indexing)
%
  bounded observer, epistemic floor, categorical observation,
  bounded phase space, cellular health, metabolic circuits,
  Wegscheider condition, holonomy, loop closure, self-consistency,
  coherence deficit, disease modelling, reference-free diagnostics,
  no-template theorem, local invisibility, fixed-point dynamics,
  Banach contraction, sparse optimisation, computational systems biology,
  mathematical biology, philosophy of measurement


=====================================================================
BACK COVER TEXT
(approx. 1,760 characters; paste as the back-cover blurb)
=====================================================================
%
 In 1568, Pieter Bruegel painted six blind men falling into a ditch --
and gave each of them a different, clinically diagnosable disease of
the eye. Not blindness the abstraction. Six men, each blind in his own
irreducible way.
%
This monograph takes that painting seriously as a theorem. From two
axioms -- that any bounded system occupies a phase space of finite
measure, and that any finite observer partitions it into finitely many
cells -- it derives a single, uncomfortable conclusion: every finite
observer has a strictly positive epistemic floor, a minimum error no
improvement in resolution can erase. The truth about a state, for such
an observer, is never a point. It is a region. And no ensemble of the
blind corrects the blindness of any one of them.
%
Yet living tissue navigates correctly -- not because any cell sees the
destination, but because its internal loops close. Health, on this
account, is not knowledge of the correct state; no cell can hold such
knowledge. Health is self-consistency: a discrepancy that vanishes on
every closed traversal of a metabolic circuit, checkable entirely from
the inside, without any external template.
%
From this one geometric fact the work develops a full structure -- the
No-Template and Local-Invisibility theorems, disease as coherence
deficit, an exact criterion for therapeutic intervention, and a picture
of the body as a system that enacts its destination without ever seeing
it. Every theorem is stated with hypothesis and conclusion; every claim
is proved or explicitly marked a conjecture.
%
It begins with a painting of six blind men, and ends with a theorem
about what it means to be alive without knowing where you are going.
%
Health is not vision. Health is loop closure.


=====================================================================
BIOGRAPHY
(approx. 470 characters; paste as the author bio)
=====================================================================
%
Kundai Sachikonye is a researcher focused on mass-spectrometry-based
methods and computational systems biology. His experience in
diagnostics, therapy, and pharmaceutical production nudged him onto the
path that led to these insights. In his free time he occasionally
engages with items of cultural importance, among them the Lowland
Masters, Oskar Schlemmer, and Roger Ballen.


=====================================================================
END OF DETAILS PACKAGE
=====================================================================
